En las fórmulas siguientes, $n\geq1$. Los coeficientes no indicados son cero. La suma coincide con la extensión periódica en sus puntos de continuidad.

\begin{enumerate}[label=(\alph*),ref=\alph*,leftmargin=*]
\item Con $k_n=n\pi/2$,
\[
C=\frac13,\qquad a_n=\frac{2(\sin k_n-k_n\cos k_n)}{k_n^3},\qquad b_n=0.
\]
La suma vale $0$ en $x=\pm1$ y $x=\pm2$, y en sus trasladados de período $4$.
\item Con $k_n=n\pi/3$,
\[
C=0,\qquad a_n=0,\qquad b_n=\frac23\left[\frac{\cos k_n-3(-1)^n}{k_n}-\frac{\sin k_n}{k_n^2}\right].
\]
La suma vale $-1/2$ en $x=-1$, $1/2$ en $x=1$ y $0$ en $x=\pm3$, con periodicidad $6$.
\item Con $k_n=n\pi/2$,
\[
C=\frac{17}{16},\qquad
a_n=\frac{(-1)^n-\cos(k_n/2)}{k_n^2},\qquad
b_n=\frac{(-1)^n}{2k_n}-\frac{\sin(k_n/2)}{k_n^2}.
\]
La suma vale $0$ en el vértice $x=1/2$ y $2$ en $x=\pm2$, con periodicidad $4$.
\item
\[
C=0,\qquad a_n=0,\qquad b_n=\frac{12(-1)^{n+1}}{n^3\pi^3}.
\]
La suma vale $0$ en $x=\pm1$ y sus trasladados de período $2$.
\item $C=1$, $a_2=1/2$, $a_4=-1/2$, $b_1=b_3=1$; todos los demás coeficientes son cero. La serie es finita y vale $1$ en $x=\pm\pi$.
\item $C=1$, $b_n=0$ y
\[
a_1=-\frac12,\qquad a_n=\frac{2(-1)^n}{1-n^2}\quad(n\geq2).
\]
La suma vale $0$ en $x=\pm\pi$ y sus trasladados de período $2\pi$.
\end{enumerate}
